\section{不可约表示是基本构建块}
\label{sec:04-Discrete-Mathematics/09-Representation-Theory/02}

\subsection{一个具体困境：六个置换，无穷多种矩阵写法}

取 \(S_3\)——三个对象的全体置换，一共六个元素。要把它``表示''成矩阵，最简单的想法是给每个置换分配一个 \(n\times n\) 矩阵，让矩阵乘法忠实反映置换的复合。可矩阵的大小和模样有无限种选择：你可以让每个置换对应 \(3\times3\) 的置换矩阵，也可以嵌入到 \(100\times100\) 的分块对角阵里塞满零，甚至可以把同一个 \(S_3\) 行动在不同维数的空间上写出表面完全不同的矩阵。问题是：这些看似不同的表示之间有没有公共的``原子零件''？能不能把所有表示拆成一批不能再拆的基本款，然后只研究基本款就够了？

先从手头最能把握的例子下手。\(S_3\) 的置换表示：三个基向量 \(e_1,e_2,e_3\)，每个置换把它俩重新排列——

\[
\rho((12))=\begin{pmatrix}0&1&0\\1&0&0\\0&0&1\end{pmatrix},\qquad
\rho((123))=\begin{pmatrix}0&0&1\\1&0&0\\0&1&0\end{pmatrix}.
\]

一眼能看出一个不变子空间：\((1,1,1)^{\trans}\) 方向被任何置换固定，因为它对所有坐标一视同仁。这意味着整个 3 维表示可以``切下一块''1 维的，剩下一个 2 维的补块。切下来的 1 维块是平凡表示——每个群元都作用于恒等矩阵 \(I_1\)。剩下的 2 维块还能再切吗？

\subsection{尝试工具：找不变子空间，切，再切}

一般的策略是：给定表示 \(\rho:G\to GL(V)\)，找到真不变子空间 \(W\subset V\)（对所有 \(g\)，\(\rho(g)W\subseteq W\)），把 \(V\) 限制到 \(W\) 上得到一个子表示；然后想办法在商空间 \(V/W\) 上也诱导一个表示。往复操作，直到每个分块里再也找不出更小的不变子空间——这就是不可约表示：非零，且除了 \(\{0\}\) 和自身以外没有别的不变子空间。

碰到的问题是：把 \(W\) 切下来之后，剩下的那个``补块''是不是也是个不变子空间？在 \(S_3\) 的置换表示中，\((1,1,1)^{\trans}\) 的补空间——向量分量之和为零的子空间——碰巧也是不变的。但碰巧是靠不住的。

\subsection{碰壁：不是所有补空间都不变}

换一个群试试。考虑整数加群 \(\Z\) 的 2 维表示：

\[
\rho(n)=\begin{pmatrix}1&n\\0&1\end{pmatrix}.
\]

取 \(W=\{(x,0)^{\trans}:x\in\C\}\)——这是不变子空间，因为 \(\rho(n)\) 只把第一分量保持不变、第二分量往上加东西但从不把第二分量混进第一分量。但你能找到 \(W\) 的一个不变补空间吗？任何补空间必须由形如 \((a,1)^{\trans}\) 的向量张成（第二分量非零才能补全 2 维）。施加 \(\rho(n)\)：

\[
\rho(n)\begin{pmatrix}a\\1\end{pmatrix}=\begin{pmatrix}a+n\\1\end{pmatrix}.
\]

第二个分量确实不变，可第一个分量在跑——这个向量被推出了补空间。没有不变补。\(\Z\) 是无限群，有限性在这里是关键分水岭。

这就卡住了：有些表示似乎拆不干净。什么时候能拆干净？

\subsection{粗糙直觉：把投影平均掉}

回到有限群。取 \(W\) 的任意线性补空间 \(U\)（纯线性代数保证存在，只是未必 \(G\)-不变），构造投影算子 \(P:V\to W\) 沿着 \(U\) 的方向。\(P\) 满足 \(P^2=P\)，像空间是 \(W\)，核是 \(U\)。但 \(P\) 一般不跟群作用交换：\(\rho(g)P\rho(g)^{-1}\) 是一个新的投影算子，像空间仍然是 \(W\)（因为 \(W\) 不变），核变成了 \(\rho(g)U\)。

现在发挥有限群的优势：把所有共轭投影取平均。定义

\[
\bar P=\frac{1}{|G|}\sum_{g\in G}\rho(g)P\rho(g)^{-1}.
\]

这是有限个算子的平均，良定义。它的像空间还是 \(W\)，但关键在——对任意 \(h\in G\)，

\[
\rho(h)\bar P\rho(h)^{-1}
=\frac{1}{|G|}\sum_{g\in G}\rho(hg)P\rho(hg)^{-1}
=\bar P,
\]

求和只是把群元重新排列了一遍。所以 \(\bar P\) 与所有 \(\rho(h)\) 交换，它的核 \(\ker\bar P\) 就是一个不变补空间。把 \(P\) 粗暴地取平均，就把一个普通投影强制改造成了 \(G\)-不变的投影。有限性使平均是有限和，分母 \(|G|\) 在底域中可逆——复数域自动满足，一般域要求特征不整除 \(|G|\)。

\subsection{精确陈述：Maschke 定理与完全可约性}

\textbf{定理（Maschke）.} 设 \(G\) 是有限群，\(\rho:G\to GL(V)\) 是复数域上的有限维表示。若 \(W\subseteq V\) 是不变子空间，则存在不变补空间 \(W'\) 使

\[
V=W\oplus W'.
\]

反复使用这个定理：把整个空间切成不可约不变子空间的直和。这就是完全可约性——复数域上有限群的每个有限维表示都是不可约表示的直和。各个不可约同构类型出现的重数唯一确定（不计次序和同构），但具体的分解子空间不唯一：同一个不可约成分可以``嵌''在不同的子空间位置里。

把``研究所有表示''化简为``研究不可约表示''——剩下的只是搞清楚一个给定表示由哪些不可约块拼成，每块用了几次。

\subsection{数值落地：不可约表示的个数与维数}

有限群 \(G\) 到底有多少个不等价的不可约复表示？答案漂亮得不像话：恰好等于 \(G\) 的共轭类个数。维数满足平方和公式：

\[
\sum_i d_i^2=|G|,
\]

其中 \(d_i\) 是第 \(i\) 个不可约表示的维数。

\(S_3\) 有三个共轭类：\(\{e\}\), \(\{(12),(13),(23)\}\), \(\{(123),(132)\}\)。所以恰好三个不可约表示：1 维平凡表示（每个群元 \(\mapsto 1\)），1 维符号表示（偶置换 \(\mapsto 1\)，奇置换 \(\mapsto -1\)），2 维标准表示（把 \(S_3\) 看作等边三角形的对称群作用于平面）。维数平方和：\(1^2+1^2+2^2=6=|S_3|\)。

\(D_4\)（正方形对称群，8 阶）有 5 个共轭类，5 个不可约表示：四个 1 维，一个 2 维。平方和 \(1+1+1+1+4=8\)。

背后的结构原因：群代数 \(\C[G]\) 作为 \(G\)-模分解为矩阵代数的直和，每个矩阵代数对应一个不可约表示，维数 \(d_i\) 的平方正是该矩阵块的维数。

\subsection{特征标：用迹查表代替找矩阵}

表示 \(\rho\) 的特征标定义为一个函数

\[
\chi_\rho(g)=\trace\rho(g).
\]

迹不依赖基的选择——这是特征标成为表示不变量的根本理由。而且它是类函数：共轭元素有相同迹，\(\chi(hgh^{-1})=\chi(g)\)。

三件关键事实：

\begin{itemize}
\tightlist
\item
  复数域上，不可约表示由其特征标唯一确定：\(\chi_\rho=\chi_{\rho'}\) 当且仅当 \(\rho\cong\rho'\)；
\item
  特征标表是方阵：行对应不可约表示，列对应共轭类，交点填特征标值。\(S_3\) 的表是 \(3\times3\)，\(D_4\) 的表是 \(5\times5\)；
\item
  正交关系：不同不可约表示的特征标关于群平均正交，
  \[
  \frac1{|G|}\sum_{g\in G}\overline{\chi_i(g)}\chi_j(g)=\delta_{ij}.
  \]
\end{itemize}

操作上的意义：给定一个表示，算出每个群元上的迹，得到一条类函数；把它与特征标表的各行做内积，各不可约成分的重数就自动浮出来。不需要在矩阵层面肉眼找不变子空间——特征标把分解变成了查表和算内积。

\subsection{Fourier 分析是表示论的特例}

有限群 \(G\) 上的复值函数全体 \(\C[G]\)，以群的平移作为 \(G\) 的作用，是一个表示——正则表示。Maschke 定理把它完全分解：

\[
\C[G]\cong\bigoplus_i d_iV_i,
\]

第 \(i\) 个不可约表示 \(V_i\) 出现 \(d_i\) 次。任何函数 \(f:G\to\C\) 都可以展开为不可约表示矩阵元素的线性组合——这就是有限群上的 Fourier 变换。

取 \(G=\Z/n\Z\)（循环群）。所有不可约表示都是 1 维：\(\rho_k(j)=e^{2\pi ijk/n}\)，\(k=0,\ldots,n-1\)。特征标就是这些指数函数自己。Fourier 展开恰好是离散 Fourier 变换（DFT），\(e^{2\pi ijk/n}\) 既是不可约表示的特征标，也是 DFT 的基函数。\hyperref[sec:08-Numerical-Analysis/08-Fourier-and-Spectral-Methods/03]{``DFT 是变换，FFT 是算法''}里的正交基，不是随便挑的——它们是循环群仅有的不可约表示的矩阵元素。

对一般有限群，Fourier 变换把卷积变成逐不可约块的矩阵乘法。快速 Fourier 变换在非交换群上的推广（如 Clausen--Baum 和 Diaconis--Rockmore 的工作），正是靠群代数的矩阵分解来组织计算。

\subsection{等变神经网络：不可约块决定能搭什么桥}

等变神经网络要求层映射与群作用交换——旋转输入，输出跟着旋转同样的方式。表示论把这种约束翻译为代数条件：等变线性映射的空间完全由输入和输出表示的不可约分解决定。

具体写出来：输入表示分解为 \(\bigoplus_i m_iV_i\)（\(m_i\) 是第 \(i\) 个不可约成分的重数），输出表示分解为 \(\bigoplus_j n_jV_j\)。层映射 \(T\) 满足 \(T\rho_{\text{in}}(g)=\rho_{\text{out}}(g)T\) 对所有 \(g\) 成立，则 \(T\) 的自由度是

\[
\sum_{i,j}m_in_j\dim\operatorname{Hom}_G(V_i,V_j).
\]

Schur 引理出场：对复数域上的不可约表示，\(\operatorname{Hom}_G(V_i,V_j)\) 在 \(i=j\) 时是 1 维（只有标量乘法），在 \(i\ne j\) 时是 0 维（零映射是仅有的等变映射）。结果是：等变层只能在相同的不可约成分之间传递信息，不同成分之间不存在非零的等变线性映射。

这就是参数共享的代数起源。不是手动设计哪些权重该共享——不可约分解自动告诉你哪些连接必须为零，哪些通道可以承载独立的标量参数。\hyperref[sec:04-Discrete-Mathematics/06-Groups-and-Symmetry/04]{``不变函数与等变映射：对称性怎样进入计算''}中的权重绑定模式，底下的硬理由就在这里。

\subsection{落点与边界}

不可约表示是群论中``原子''概念的精确化身。Maschke 定理保证有限群在复数域上的表示总能拆成这些原子；特征标表把分解问题变成内积计算。Fourier 分析是在循环群上展开这个框架，而等变神经网络则是在一般有限群上利用 Schur 引理收紧参数空间。

但原子图景有明确的边界。无限群的表示一般不完全可约——\(\Z\) 的那个 2 维表示就是永远的警示。即使对有限群，换到底域特征整除 \(|G|\) 的 modular 表示论中，Maschke 定理失效，不可约表示的分类远比复数域上复杂。复数域上的干净图景不是表示论的终点，它是通向更深水域的码头。
